\section{What the exam is and where the points are}

The Putnam is one Saturday, and from 2026 it has a new shape. Four sessions of 90 minutes, A, B, C and D, with three problems each: A1 to A3, B1 to B3, C1 to C3 and D1 to D3. Until 2025 it was two sessions of three hours with six problems each, A1 to A6 and B1 to B6, and every past paper named in this document has that older shape. The total is unchanged: twelve problems, each marked out of 10, so the paper is still out of 120. Every site sits each session at the same moment. In the Eastern time zone the sessions start at 11:00, 12:45, 16:00 and 17:45, so there are fifteen minutes between A and B, an hour and three quarters between B and C, fifteen minutes between C and D, and the day ends at 19:15. The source is the MAA's policy document for the 2026 competition, in the version updated 2026-08-31.

What the policy does not say is how difficulty is spread. In the old papers the problems of a session ran roughly from most approachable to hardest, with A1 and B1 meant for most contestants and A6 and B6 meant to be solved by very few people in the country. That each new session runs from an approachable first problem to a hard third is likely and not stated, and nobody outside the committee has seen a paper in the new shape. Treat the order as a hint and the reading of all three problems as the evidence.

\subsection*{How a problem is scored}

The scoring is close to binary. Almost every score given is 0, 1, 9 or 10. The reasons are worth knowing, because they decide how to spend a session. The 2026 policy says only that each problem is scored from 0 to 10 by at least two graders working blind, and that scores are final. The pattern below is the long practice of the old format, and nothing in the policy suggests it has changed.

\begin{itemize}[nosep,leftmargin=*]
\item 10: a complete and correct argument.
\item 9: a complete argument with a small, repairable gap. A missed equality case, a base case not written, a symbol not defined but clear from context.
\item 1 or 2: a correct idea with no proof, or a proof of a special case, or the right answer with a wrong or missing reason.
\item 0: everything else, including a page of true statements that do not reach the conclusion.
\end{itemize}

There is nearly nothing in the middle. A half-finished solution is worth about as much as a blank page. That single fact drives the strategy in this document: finishing one problem beats starting three.

\subsection*{What a score means}

The median total most years is a single digit, often 0, 1 or 2. In a typical year, one fully solved problem (10 points, or 9) puts a contestant above the median. Two fully solved problems, around 20 points, are usually in the top quarter of the roughly four thousand people who sit it. Three are near the top ten percent. These thresholds move a little from year to year with the paper's difficulty, but not by much. They come from the two-session papers. The 2026 paper is the first in the new shape and its thresholds are not known in advance. The argument does not depend on them: a finished problem is worth ten started ones under either shape.

So the realistic target for a first sitting is two problems solved and written cleanly across the four sessions, with a third attempted wherever a session leaves time. Under the old shape that target was reached almost entirely in the first two slots of each session.

\subsection*{Where the points are}

Across 2016 to 2025 there were forty problems in the slots A1, A2, B1 and B2 of the old papers, the four approachable slots. The data file behind this document tags each by the kind of object and method its statement points to; most problems carry two tags, so the forty problems carry seventy tags. The counts are in Table~\ref{tab:tags}.

\begin{table}[ht]
\centering
\small
\begin{tabular}{lr}
\toprule
Tag & Problems among A1, A2, B1, B2, 2016 to 2025 \\
\midrule
Congruences and number theory & 13 \\
Sums, integrals and limits & 10 \\
Counting and configurations & 10 \\
Polynomials and roots & 9 \\
Geometry & 6 \\
Inequalities & 5 \\
Linear algebra & 3 \\
Recurrences & 3 \\
Pigeonhole and extremal & 3 \\
Invariants and monovariants & 3 \\
Hidden symmetry & 2 \\
Induction & 1 \\
Games & 1 \\
Probability & 1 \\
Generating functions & 0 \\
Abstract algebra & 0 \\
Functional equations & 0 \\
\bottomrule
\end{tabular}
\caption{Tag counts in the four easy slots of the old papers. Forty problems, seventy tags. The tags are a model's reading of the statements, not of the solutions.}
\label{tab:tags}
\end{table}

Four topics carry more than half the tags: number theory, sums and limits, counting, and polynomials. Those four are the four weekly ladders in \secref{sec:ladders}{the four ladders}, in that order. Geometry and inequalities together take another eleven; they are second rails rather than ladders, because the seminar has thin material on geometry and because inequalities tend to appear inside a limit or a polynomial problem rather than alone.

The techniques at the bottom of the table are not useless. Pigeonhole, invariants and induction appear under other tags far more often than their own counts suggest, which is why each technique section ends with a note on where it hides. But they are not where the first half of a session goes.

The new paper has four sessions where the old one had four approachable slots, and the natural guess is that each session opens with a problem of that kind. That is a guess; the policy says nothing about it. What carries over without guessing is the table: ten years of what the committee chooses for the problems most contestants can reach.

\subsection*{What this means for the month}

Four weeks, four ladders, in the order of the table. Each ladder climbs one technique through its easy slots and then into one harder problem, so that the easy slots feel easy on the day. Two timed half-papers, on 2026-11-28 and 2026-12-02, each cut into two 90-minute sittings of three problems, rehearse the new shape. The target on 2026-12-05 is two clean solutions across the four sessions, and the habit of reading all three problems of a session before touching any of them.

